\chapter{Material complementario}
\label{app:material-complementario}

\section{Esquema del conjunto de resultados}

El archivo tabular de resultados contiene una fila por imagen y veinticuatro columnas,
organizadas en cinco grupos. La Tabla~\ref{tab:esquema} las describe.

\begin{table}[H]
  \centering
  \small
  \caption{Esquema del conjunto de resultados.}
  \label{tab:esquema}
  \begin{tabular}{llp{6.0cm}}
    \toprule
    \textbf{Grupo} & \textbf{Columna} & \textbf{Descripción} \\
    \midrule
    Identificación & \texttt{image\_id}   & Identificador de la escena. \\
                   & \texttt{rover}       & Misión de procedencia. \\
                   & \texttt{camera}      & Cámara de procedencia. \\
                   & \texttt{eye}         & Ojo del par estéreo (L o R). \\
                   & \texttt{sol}         & Día marciano, si está disponible. \\
    \midrule
    Cobertura (E1) & \texttt{rock\_coverage\_pct}  & Roca sobre píxeles etiquetados. \\
                   & \texttt{coverage\_total\_pct} & Roca sobre la imagen completa. \\
                   & \texttt{frac\_valid}          & Fracción de escena etiquetada. \\
    \midrule
    Composición    & \texttt{pct\_soil}       & Proporción de suelo. \\
                   & \texttt{pct\_bedrock}    & Proporción de lecho rocoso. \\
                   & \texttt{pct\_sand}       & Proporción de arena. \\
                   & \texttt{pct\_bigrock}    & Proporción de roca grande. \\
                   & \texttt{dominant\_class} & Clase mayoritaria. \\
                   & \texttt{scene\_type}     & Tipología de escena. \\
    \midrule
    Conteo (E2)    & \texttt{n\_bigrock}          & Píxeles de roca grande. \\
                   & \texttt{n\_rocks}            & Rocas tras división de aguas y filtros. \\
                   & \texttt{n\_raw\_components}  & Componentes conectadas previas. \\
                   & \texttt{largest\_rock\_pct}  & Tamaño de la roca mayor. \\
                   & \texttt{mean\_rock\_area\_px}& Área media de las rocas. \\
                   & \texttt{n\_small}            & Rocas pequeñas. \\
                   & \texttt{n\_medium}           & Rocas medianas. \\
                   & \texttt{n\_large}            & Rocas grandes. \\
    \midrule
    Calidad        & \texttt{mean\_solidity} & Solidez media de las regiones aceptadas. \\
                   & \texttt{quality\_flag}  & Bandera de aptitud de la escena. \\
    \bottomrule
  \end{tabular}
\end{table}

\section{Casos adicionales del procedimiento}

Las figuras de esta sección muestran el procedimiento sobre escenas de tipos de terreno
distintos, y sirven de complemento a la Figura~\ref{fig:paso-a-paso}.

\begin{figure}[H]
  \centering
  \includegraphics[width=\textwidth]{Figura_A1_procedimiento_anexo_1.png}
  \caption{Procedimiento sobre una escena con una roca aislada.}
  \label{fig:anexo1}
\end{figure}

\begin{figure}[H]
  \centering
  \includegraphics[width=\textwidth]{Figura_A2_procedimiento_anexo_2.png}
  \caption{Procedimiento sobre una escena con varios bloques próximos.}
  \label{fig:anexo2}
\end{figure}

\begin{figure}[H]
  \centering
  \includegraphics[width=\textwidth]{Figura_A3_procedimiento_anexo_3.png}
  \caption{Procedimiento sobre una escena de terreno mixto.}
  \label{fig:anexo3}
\end{figure}

\begin{figure}[H]
  \centering
  \includegraphics[width=\textwidth]{Figura_A4_procedimiento_anexo_4.png}
  \caption{Procedimiento sobre una escena dominada por afloramiento.}
  \label{fig:anexo4}
\end{figure}

\section{Modos de error descartados por los filtros}

Las dos figuras siguientes documentan los casos que el
Tabla~\ref{tab:diagnostico} atribuye a los umbrales del procedimiento, y respaldan la
afirmación de que en esas escenas los filtros aciertan.

\begin{figure}[H]
  \centering
  \includegraphics[width=\textwidth]{diag_bajo_umbral_area.png}
  \caption[Región por debajo del área mínima]{Escena con una única región de roca grande
    de doce píxeles, por debajo del umbral de \num{524}. El conteo devuelve cero
    correctamente.}
  \label{fig:diag-area}
\end{figure}

\begin{figure}[H]
  \centering
  \includegraphics[width=\textwidth]{diag_filtro_aspecto.png}
  \caption[Región descartada por relación de aspecto]{Escena en la que la anotación de
    roca grande es una banda alargada trazada sobre el horizonte, descartada por relación
    de aspecto.}
  \label{fig:diag-aspecto}
\end{figure}

\begin{figure}[H]
  \centering
  \includegraphics[width=\textwidth]{diag_sobresegmentacion.png}
  \caption[Anotación en contorno hueco]{Caso atípico de anotación trazada como contorno
    hueco. La división de aguas fragmenta el anillo resultante. Se midió que este patrón
    afecta a cerca del \SI{0.2}{\percent} de las escenas con roca grande.}
  \label{fig:diag-sobreseg}
\end{figure}

\section{Documentación técnica de la extensión aplicada}
\label{app:extension}

Esta sección documenta los componentes descritos en la
subsección~\ref{sec:extension}.

\subsection*{Capa de avisos}

Se implementa como un módulo independiente que opera sobre el archivo de resultados y le
añade cuatro columnas: las claves de los avisos activados, su número, la severidad máxima
y el nivel de riesgo resultante. Cada regla se declara como un objeto con su clave, su
título legible, su severidad, la condición —una función sobre la fila de resultados—, el
criterio en texto y su motivación. Esa estructura permite añadir, retirar o reajustar
reglas sin tocar la lógica de evaluación, y garantiza que todo umbral quede acompañado de
su justificación en el mismo lugar del código.

La evaluación no reprocesa máscaras: recorre el archivo de resultados, de modo que
recalcular los avisos sobre las \num{16064} escenas es inmediato. Las salidas son la tabla
completa con las columnas añadidas, un catálogo de las reglas con sus umbrales y un
resumen agregado por nivel y por tipo de aviso.

\subsection*{Aplicación de consulta}

Aplicación de escritorio en Python construida sobre una biblioteca de componentes gráficos
modernos, con la paleta institucional de la agencia responsable del conjunto de datos. La
ventana se organiza en una barra lateral de navegación y cuatro vistas:

\begin{description}
  \item[Resumen] Estadísticos descriptivos del subconjunto y composición del terreno.
  \item[Avisos] Distribución por nivel de riesgo y por tipo de regla, con el detalle del
    criterio que motivó cada aviso.
  \item[Explorador] Para la escena seleccionada, la imagen original, la anotación del
    terreno y las rocas detectadas por el procedimiento, en paneles contiguos.
  \item[Geología] Presencia de los rasgos de la taxonomía extendida en el conjunto de la
    misión posterior, incluidas las vetas.
\end{description}

Las figuras se generan en el momento a partir de los mismos archivos de resultados que
respaldan este documento, de modo que la aplicación no mantiene una copia propia de los
datos y cualquier reejecución del procedimiento se refleja en ella sin pasos adicionales.

\subsection*{Exploración de rasgos geológicos}

El conjunto de datos incluye, para la misión posterior a la estudiada, una taxonomía
geológica más detallada que la escala de navegación: distingue subtipos de lecho rocoso,
roca suelta, guijarros y vetas. Se recorrieron sus \num{8497} máscaras registrando la
presencia y extensión de cada rasgo. Las vetas aparecen en \num{294} escenas, el
\SI{3.5}{\percent}, de las cuales \num{150} presentan una geometría compatible con una
veta genuina —alargada y de área reducida— y \num{69} resultan dudosas.

Ese recorrido tiene un interés que excede la exploración de vetas y se recoge entre las
recomendaciones del Capítulo~\ref{cap:conclusiones}: la clase de \textit{roca suelta} de
esa taxonomía aparece en una proporción muy superior a la de \textit{roca grande} en la
escala de navegación, lo que la convierte en la vía natural para sortear el techo del
conteo documentado en la subsección~\ref{sec:diagnostico}.

\section{Recursos reproducibles}

Los recursos que respaldan el trabajo se organizan como sigue.

\begin{description}

  \item[Módulos de cálculo] Implementan la lectura de máscaras, el cálculo de cobertura,
    el conteo de rocas, los descriptores de composición y geometría, y la orquestación
    que produce una fila de resultados por imagen. Las funciones reciben arreglos y
    devuelven arreglos o números, sin efectos secundarios, y todos los umbrales se pasan
    como argumentos con valores por defecto documentados.

  \item[Guiones de ejecución] Automatizan el recorrido del subconjunto, la generación de
    las figuras del documento, el diagnóstico de los modos de error y las comparaciones
    con los métodos de aprendizaje automático.

  \item[Conjunto de resultados] Archivo tabular con los veinticuatro indicadores para las
    \num{16064} escenas, según el esquema de la Tabla~\ref{tab:esquema}.

  \item[Especificación del entorno] Archivo de definición del entorno con las
    dependencias, complementado por el registro de versiones exactas del
    Tabla~\ref{tab:versiones}.

\end{description}

El conjunto de datos AI4Mars no se distribuye con estos recursos, dado su volumen. El
código lo referencia mediante una ruta configurable, de modo que debe descargarse de la
fuente original \citep{swan2022ai4marsdata}.
